\section{Results III: Where No Candidate Budget Suffices}
\label{sec:results3}

\textbf{Mode-probability flattening.}
\label{sec:flatten}
On a bounce variant with $K{=}2$ modes of probability $0.8$ and $0.2$ (realized frequency $0.787$ on the evaluation contexts), the share of the frequent mode among 64 samples per context falls with $m$ as the stationarity condition of \citet{bennett2026freedom} prescribes: $0.82$, $0.65$ and $0.57$ for $m = 2$, $4$ and $16$, against theoretical values of $0.79$, $0.61$ and $0.52$ (Figure~\ref{fig:knob}d; Table~\ref{tab:flatten}, Appendix~\ref{app:more}). Every method, including FM-50 ($0.83$), measures about $0.04$ above its reference value, a common offset of the readout-based assignment; relative to FM-50 the measured flattening is $-0.18$ and $-0.26$, against theoretical $-0.18$ and $-0.27$. The trade-off is real: $m{=}2$ is calibrated but memorizes (index $0.93$) and covers only $0.75$; $m{=}16$ covers $0.99$ but reports the modes as nearly equiprobable; FM-50 does both; and the energy score does neither (share $0.89$, coverage $0.58$). Inverting the stationarity condition to recover $p$ is ill-conditioned for large $m$ (errors in the log-odds are amplified about $60\times$ at $m{=}16$), so a large budget is a coverage device, not a probability estimate; for equiprobable modes, the calibration is exact (Section~\ref{sec:planning}).


\textbf{Under-dispersion on real robot video (BAIR).}
\label{sec:bair}
BAIR \citep{ebert2017bair} has no enumerable modes; its uncertainty is a continuous, high-dimensional variation of the arm and the objects. \imlewm{} is the most accurate model there (best-of-16 token error $0.527$ at $h{=}16$ against $0.843$ for FM-50; manifold distance $1.30$, like regression, against $1.78$) but loses clearly on distribution metrics (Fr\'echet distance $2.09$ vs.\ $1.38$ and kNN recall $0.15$ vs.\ $0.84$ for FM-50; Table~\ref{tab:bair}, Appendix~\ref{app:more}). The samples are under-dispersed: their spread-to-error ratio at $h{=}1$ is $0.22$ at $m{=}4$ and $0.45$ at $m{=}16$, the same on training and evaluation contexts, so this is not memorization. Raising the budget to $m{=}64$ moves the ratio only to $0.54$ at $2.4\times$ the training time, and matching in a 512-dimensional pooled space lowers it to $0.31$--$0.38$; no one-step \imlewm{} variant closes the gap. Proposition~\ref{prop:shrink} reproduces the growth of the spread with $m$ using one fitted parameter ($\deff \approx 27$, fitted at $m{=}16$, predicts $0.27$ and $0.58$ at $m = 4$ and $64$ against measured values of $0.22$ and $0.54$; Figure~\ref{fig:bair}a, Appendix~\ref{app:more}), but our recorded attempt to remove that parameter failed: the participation ratio of the regression residuals is about $1{,}030$, far outside a factor of two of $27$, and would predict a spread five times smaller than observed. An exploratory analysis designed after this failure suggests why: per context, the samples occupy only $1.2$, $1.8$ and $2.5$ effective directions at $m = 4$, $16$ and $64$; these directions differ across contexts; and they align with the true residual of their own context ten times better than with that of another context, while capturing only $9$--$21\%$ of it (Figure~\ref{fig:bair}b,c). The contraction is therefore a collapse onto a few context-specific directions whose number grows slowly with $m$, the shape of reverse water-filling in rate--distortion theory; with a single true future per context we could not test this quantitatively.
